% answers_short.tex —— 2020 年·新高考Ⅰ卷（山东） · 简短答案速查
% 由 tools/make_short_answers.py 自 answers.tex 的 \gksolhead/\gkanswer 提取，
% 并与 2020/2020_数学真题_新高考1_答案要点.md 逐题比对后生成，勿手改（改源 MD 后重跑）。

\begin{gknotice}
\textbf{使用说明}：本节只给最终答案便于快速核对，完整推导见解析版。答案为编写者独立推导，\textbf{非官方标准答案}。
\end{gknotice}

\gkpart{单选题}
\begin{center}\begin{tabular}{ccccccccc}
\toprule
题号 & 1 & 2 & 3 & 4 & 5 & 6 & 7 & 8 \\
\midrule
答案 & C & D & C & B & C & B & A & D \\
\bottomrule
\end{tabular}\end{center}

\gkpart{多选题}
\begin{center}\begin{tabular}{ccccc}
\toprule
题号 & 9 & 10 & 11 & 12 \\
\midrule
答案 & ACD & BC & ABD & AC \\
\bottomrule
\end{tabular}\end{center}

\gkpart{填空题}
\begin{itemize}[leftmargin=2.6em]
  \item[\textbf{13.}] $\dfrac{16}3$。焦点 $F(1,0)$，直线 $y=\sqrt3(x-1)$，代入得 $3x^2-10x+3=0$，$x_1+x_2=\dfrac{10}3$，故 $|AB|=x_1+x_2+p=\dfrac{16}3$。
  \item[\textbf{14.}] $3n^2-2n$。公共项即除以 $6$ 余 $1$ 的正整数，$a_n=6n-5$，$S_n=\dfrac{n(1+6n-5)}2=3n^2-2n$。
  \item[\textbf{15.}] $4+\dfrac{5\pi}2$。建系得 $E(0,0)$、$F(12,0)$、$G(12,2)$、$D(0,2)$、$A(7,7)$；由 $OA\perp AG$ 与 $\tan\angle ODC=\dfrac35$ 得 $O(5,5)$、$R=2\sqrt2$、$B\left(5-2\sqrt2,5\right)$、$H(9,5)$，$\angle AOB=135^\circ$；阴影面积 $=\dfrac{135}{360}\pi R^2+\dfrac12\cdot OA\cdot AH-\dfrac12\pi\cdot1^2=3\pi+4-\dfrac\pi2$。
  \item[\textbf{16.}] $\dfrac{\sqrt2\pi}2$。建系得 $D_1$ 到平面 $BCC_1B_1$ 的距离为 $\sqrt3$，交线圆半径为 $\sqrt{5-3}=\sqrt2$，圆心为 $B_1C_1$ 的中点；该圆被侧面两条竖直棱截得圆心角 $90^\circ$ 的弧，长为 $\sqrt2\cdot\dfrac\pi2=\dfrac{\sqrt2\pi}2$。
\end{itemize}

\gkpart{解答题}
\begin{itemize}[leftmargin=2.6em]
  \item[\textbf{17.}] 选① $c=1$；选② $c=2\sqrt3$；选③ 三角形不存在。公共条件给出 $a=\sqrt3b$、$c=b$：选①由 $ac=\sqrt3$ 得 $b=1$、$c=1$；选②由 $c\sin A=3$ 得 $a=6$、$b=c=2\sqrt3$；选③要求 $c=\sqrt3b$ 与 $c=b$ 同时成立，只能 $b=0$，矛盾。
  \item[\textbf{18.}] (1) $a_n=2^n$。由 $a_3=8$ 得 $a_1=\dfrac8{q^2}$，代入 $a_2+a_4=20$ 得 $2q^2-5q+2=0$，$q=2$（舍 $\dfrac12$），$a_1=2$；(2) $480$。$b_m=\left\lfloor\log_2m\right\rfloor$，按 $m$ 所在区间 $[2^k,2^{k+1}-1]$ 分段求和：$0\cdot1+1\cdot2+2\cdot4+3\cdot8+4\cdot16+5\cdot32+6\cdot37=480$
  \item[\textbf{19.}] (1) $0.64$。落在「PM2.5 不超过 $75$ 且 $\mathrm{SO_2}$ 不超过 $150$」的天数为 $32+18+6+8=64$，频率为 $\dfrac{64}{100}=0.64$；(2) $a=64$，$b=16$，$c=10$，$d=10$（列联表见解析正文）。行合计 $80$、$20$，列合计 $74$、$26$；(3) 有 $99\%$ 的把握认为二者有关。$K^2=\dfrac{100(64\times10-16\times10)^2}{80\times20\times74\times26}\approx7.484>6.635$，但 $<10.828$
  \item[\textbf{20.}] (1) 证明。由 $BC\parallel AD$ 得 $BC\parallel$ 平面 $PAD$；又 $BC\subset$ 平面 $PBC$ 且两平面交于 $l$，故 $BC\parallel l$。由 $PD\perp$ 底面得 $PD\perp BC$，正方形给出 $BC\perp CD$，于是 $BC\perp$ 平面 $PDC$，从而 $l\perp$ 平面 $PDC$；(2) $\dfrac{\sqrt6}3$。以 $D$ 为原点建系，$l$ 过 $P(0,0,1)$ 且平行 $BC$，设 $Q(t,0,1)$，平面 $QCD$ 的法向量为 $(1,0,-t)$，$\sin\theta=\dfrac{|1+t|}{\sqrt3\sqrt{1+t^2}}\leqslant\dfrac{\sqrt2}{\sqrt3}=\dfrac{\sqrt6}3$，$t=1$ 时取到
  \item[\textbf{21.}] (1) $\dfrac{2}{\mathrm e-1}$。$a=\mathrm e$ 时 $f(x)=\mathrm e^x-\ln x+1$，$f(1)=\mathrm e+1$，$f'(1)=\mathrm e-1$，切线 $y=(\mathrm e-1)x+2$ 的两截距为 $2$ 与 $-\dfrac2{\mathrm e-1}$，面积 $\dfrac12\cdot2\cdot\dfrac2{\mathrm e-1}=\dfrac2{\mathrm e-1}$；(2) $[1,+\infty)$。取 $x=1$ 得必要条件 $a+\ln a\geqslant1$，故 $a\geqslant1$；反之 $a\geqslant1$ 时 $f(x)\geqslant\mathrm e^{x-1}-\ln x\geqslant1$（后者在 $x=1$ 取最小值 $1$）
  \item[\textbf{22.}] (1) $\dfrac{x^2}6+\dfrac{y^2}3=1$。由 $e^2=\dfrac12$ 得 $a^2=2b^2$，代入 $A(2,1)$ 得 $\dfrac4{2b^2}+\dfrac1{b^2}=1$，$b^2=3$；(2) 存在定点 $Q\left(\dfrac43,\dfrac13\right)$，使 $|DQ|=\dfrac{2\sqrt2}3$。设 $MN:y=kx+m$，由 $AM\perp AN$ 化简得 $m=-\dfrac{2k+1}3$，即 $MN$ 恒过 $P_0\left(\dfrac23,-\dfrac13\right)$；$AD\perp MN$ 使 $D$ 在以 $AP_0$ 为直径的圆上，取 $Q$ 为 $AP_0$ 中点即得定值 $\dfrac{|AP_0|}2=\dfrac{2\sqrt2}3$
\end{itemize}
